\section{행동모형: 위험선호, 확률가중과 3단계 비축약}

주분석은 가격벡터 사이의 정합성을 측정하지만 그 차이가 왜 생기는지는 식별하지
않는다. 이 절은 실현 착순과 추가 유지가정을 도입하여
\citet{snowbergwolfers2010}의 다중승식 식별논리를 세 번째 순위까지 확장한다.
\citet{koivurantakorhonen2019}의 완전한 exotic 가격 분석을 상위 3위로
확장하여, 각 풀의 전체 가격벡터를 표본 밖에서 평가하는 것이 목표다.

\subsection{조건부 순위확률의 교차적합}

경주 $r$에서 $i$가 우승할 확률을 $\pi_{r,i}$, $i$가 우승한 조건 아래 $j$가
2위일 확률을 $\pi_{r,j\mid i}$, $i,j$가 1·2위인 조건 아래 $k$가 3위일 확률을
$\pi_{r,k\mid ij}$라 한다. 그러면 순서 있는 상위 3위의 결합확률은
\begin{equation}
  \pi_{r,ijk}
  =\pi_{r,i}\pi_{r,j\mid i}\pi_{r,k\mid ij}
  \label{eq:rank-probability-factorization}
\end{equation}
이다.

이 확률은 가격 적합도를 높이기 위해 같은 경주의 실현결과에 맞추지 않는다.
연도를 fold로 둔 교차적합을 사용하여 다른 연도의 착순자료에서 추정한 모형으로
해당 연도의 모든 조합확률을 예측한다. 첫 번째 명세는 단승 확률을 순차적으로
재조정하는 Harville/Plackett--Luce 기준이고, 두 번째는 단승 가격벡터,
출전두수와 사전 고정한 경주 특성으로 2·3위 조건부분포를 직접 추정하는 유연한
순위모형이다. 구조적 가격모형을 비교하기 전에 각 단계의 합계 제약, 로그점수와
calibration을 검증한다. 이 확률모형이 표본 밖 착순을 설명하지 못하면 행동모형의
해석을 보류한다.

\subsection{단순복권에서 추정하는 두 함수}

총지급배율 $D$에 한 단위를 걸면 적중 시 순이익은 $D-1$, 실패 시 손실은 1이다.
초기부를 $W$라 할 때 위험선호 모형 M-U는 객관확률 $\pi_e$인 사건 $e$의
균형배당이
\begin{equation}
  \pi_e u(W+D_e-1)+(1-\pi_e)u(W-1)=u(W)
  \label{eq:risk-utility-pricing}
\end{equation}
을 만족한다고 둔다. 효용함수 $u$는 단승식 훈련표본에서만 추정하고, 그 뒤에는
고정하여 다른 승식의 배당을 예측한다.

확률가중 모형은 위험중립인 대표 행위자가 객관확률 $\pi$ 대신 $w(\pi)$를
사용한다고 둔다. 복합사건을 하나의 단순복권으로 축약하는 M-R의 가격식은
\begin{equation}
  w(\pi_e)D_e=1
  \label{eq:reduced-weighting}
\end{equation}
이다. $w$도 단승식 훈련표본에서 고정한다. 단조 비모수 추정을 기본으로 하고,
\citet{prelec1998}의 함수는 문헌 비교를 위한 보조 명세로 사용한다.

M-U와 M-R을 제한 없이 각각 단승 가격에 맞추면 단순복권에서는 관측적으로
구분되지 않을 수 있다 \citep{snowbergwolfers2010}. 따라서 두 모형의 비교는
사전에 고정한 함수형과 풀 밖 예측에 조건부이다. 이 논문의 더 직접적인 식별
대상은 결합확률을 한 번 평가하는 식 \eqref{eq:reduced-weighting}과 조건부
확률을 단계마다 평가하는 비축약 모형의 차이다.

\subsection{2단계와 3단계 비축약}

쌍승식 $ij$를 1위와 조건부 2위의 두 단계로 평가하는 M-S2는
\begin{equation}
  w(\pi_{r,i})w(\pi_{r,j\mid i})D^{\Eact}_{r,ij}=1
  \label{eq:sequential-exacta}
\end{equation}
을 예측한다. 완전 축약 M-R은 같은 사건에
$w(\pi_{r,i}\pi_{r,j\mid i})$를 한 번 적용한다. 확률가중 함수가 일반적인
비선형 함수이면 두 예측은 다르다.

삼쌍승에서는 세 모형을 비교한다. M-R은 결합확률 $\pi_{r,ijk}$에 $w$를 한 번
적용한다. M-S2는 1위와 그 조건 아래 2·3위 결합사건을 분리한다. M-S3은
\begin{equation}
  w(\pi_{r,i})w(\pi_{r,j\mid i})w(\pi_{r,k\mid ij})
  D^{\T}_{r,ijk}=1
  \label{eq:sequential-trifecta}
\end{equation}
처럼 세 조건부 단계를 모두 분리한다. 이는 복합복권의 축약공리를 위반하는
특정한 사건트리 표현이며 \citep{segal1990}, M-S3의 우위는 이 표현과
양립한다는 뜻이지 개인의 인지과정을 직접 관찰했다는 뜻은 아니다.

\subsection{풀 밖 검증과 모형 판정}

행동모형은 정규화 가격만이 아니라 원시 청구권 가격 $1/D$의 수준과 형태를
예측해야 한다. 공식 공제율이 확인되면 승식별 수준을 먼저 조정한다. 공제율
정보가 불완전하면 훈련표본에서 승식별 수준상수 하나만 추정하고 검증표본에서는
고정한다. 풀마다 다른 비선형 함수를 허용하여 적합도를 기계적으로 높이지 않는다.

모형 판정은 조합별 절대가격오차, 로그가격오차와 표본 밖 우도로 한다. 함수와
초모수의 추정, 수준상수의 추정, 최종 평가에는 서로 다른 fold를 사용한다.
시간순 검증과 leave-one-pool-out 검증도 함께 보고한다. 한 풀에서 다시 맞춰야만
성공하는 모형은 공통 행동모형으로 채택하지 않는다.

순서가 없는 복승·삼복승에는 자의적인 단계 순서를 부여하지 않는다. 대신
두 가지 예측을 비교한다. 하나는 순서 없는 결합사건의 확률에 $w$를 한 번
적용한 축약가격이고, 다른 하나는 쌍승의 두 순서 또는 삼쌍승의 여섯 순서에서
나온 예측 청구권 가격을 합한 값이다. 이 비교는 exacta--quinella의 상대가격을
이용한 Snowberg and Wolfers의 검정을 상위 3위까지 확장하면서, 복승·삼복승을
순서형 모형의 외부 검증으로 남긴다.

집계가격만으로는 동일한 사람이 여러 풀에 참가했는지 알 수 없다. 따라서
순차모형이 우월해도 이를 모든 개인의 확률오인으로 해석하지 않는다. 풀별 참가자
구성과 유동성을 통제한 뒤에도 함수가 승식과 연도 사이에서 이동할 때에만
안정적인 대표 행동모형과 양립한다고 결론 내린다.
